\section{September 24, 2026}

We extend integration to complex-valued functions, explain why
integrable functions are identified almost everywhere, and introduce
the dominated convergence theorem. Throughout, $(X,\mathcal M,\mu)$ is
a measure space.

\subsection{Complex-valued integrable functions}

\begin{definition}
  A measurable function $f\colon X\to\bb C$ is \emph{integrable} if
  $\operatorname{Re}f$ and $\operatorname{Im}f$ are integrable. Define
  \[
    \int_X f\,d\mu
      =\int_X\operatorname{Re}f\,d\mu
       +i\int_X\operatorname{Im}f\,d\mu.
  \]
  Denote the space of these functions by
  $\mathcal L^1(X,\mu;\bb C)$, or simply $\mathcal L^1(X,\mu)$.
\end{definition}

The inequalities
\[
  |\operatorname{Re}f|,\ |\operatorname{Im}f|\leq|f|
  \leq|\operatorname{Re}f|+|\operatorname{Im}f|
\]
show that, just as for real-valued functions,
\[
  f\in\mathcal L^1(X,\mu)
  \quad\Longleftrightarrow\quad\int_X|f|\,d\mu<\infty.
\]
The space $\mathcal L^1(X,\mu)$ is a complex vector space, and its
integral is complex linear. To see explicitly why complex scalars are
allowed, write $f=u+iv$ and $c=a+ib$. Then
\[
  cf=(au-bv)+i(bu+av),
\]
and real linearity applied to both components gives
$\int cf\,d\mu=c\int f\,d\mu$. Additivity follows componentwise.
Also, $|af+bg|\leq|a||f|+|b||g|$ proves closure under complex linear
combinations.

\begin{example}[Lebesgue and Riemann integrability]
  \leavevmode
  \begin{enumerate}[label=(\roman*)]
    \item The characteristic functions $\mathbf 1_{\bb Q}$ and
      $\mathbf 1_C$ belong to $\mathcal L^1(\bb R,m)$ and have integral
      zero, where $C$ is the Cantor set. On $[0,1]$,
      $\mathbf 1_{\bb Q}$ is not Riemann integrable because every
      lower sum is zero and every upper sum is one.
    \item Define $f(x)=\sin x/x$ for $x\neq0$ and $f(0)=1$.
      Its improper Riemann integral over $\bb R$ converges, but
      $f\notin\mathcal L^1(\bb R,m)$. For convergence of the positive
      tail, integration by parts gives, for $B>A\geq1$,
      \[
        \int_A^B\frac{\sin x}{x}\,dx
          =\left[-\frac{\cos x}{x}\right]_A^B
             -\int_A^B\frac{\cos x}{x^2}\,dx.
      \]
      The boundary term has a limit and the remaining integral
      converges absolutely. The negative tail is the same by symmetry.
      On the other hand, for every integer $k\geq1$,
      \[
        \int_{k\pi}^{(k+1)\pi}\frac{|\sin x|}{x}\,dx
          \geq\frac1{(k+1)\pi}
            \int_{k\pi}^{(k+1)\pi}|\sin x|\,dx
          =\frac2{(k+1)\pi}.
      \]
      Summing and using divergence of the harmonic series proves
      $\int_{\bb R}|f|\,dm=+\infty$. In fact the positive and negative
      parts both have infinite integral, so the signed Lebesgue
      integral is undefined.
    \item The function $x^{-1/2}$ is Lebesgue integrable on $(0,1)$,
      with
      \[
        \int_{(0,1)}x^{-1/2}\,dm
          =\lim_{n\to\infty}\int_{1/n}^1x^{-1/2}\,dx=2.
      \]
      Here monotone convergence justifies the limit of the truncated
      integrals. The function is unbounded near zero, so it has no
      proper Riemann integral on $[0,1]$, although its improper
      Riemann integral converges.
    \item The Gaussian $e^{-x^2}$ is Lebesgue integrable on $\bb R$.
      For $|x|\geq1$, $e^{-x^2}\leq e^{-|x|}$, while on $[-1,1]$
      it is bounded by $1$. These bounds give a finite integral.
      It is also improperly Riemann integrable.
  \end{enumerate}
  In the elementary computations above, we use agreement of the
  Riemann and Lebesgue integrals for continuous functions on compact
  intervals, which follows by squeezing between their lower and
  upper step functions.
\end{example}

\Needspace{7\baselineskip}
\subsection{The integral triangle inequality}

\begin{proposition}[Triangle inequality and its equality case]
  \label{prop:integral-triangle-sep-24}
  If $f\in\mathcal L^1(X,\mu)$, then
  \[
    \left|\int_X f\,d\mu\right|\leq\int_X|f|\,d\mu.
  \]
  Equality holds if and only if there is a constant $c\in\bb C$ with
  $|c|=1$ such that $f=c|f|$ a.e.
\end{proposition}

\begin{proof}
  Put $z=\int_X f\,d\mu$. If $z=0$, the inequality is immediate.
  In this case equality holds precisely when $\int|f|\,d\mu=0$,
  or equivalently $f=0$ a.e. by Lemma~\ref{lem:zero-integral-sep-22}.
  Such a function satisfies $f=c|f|$ a.e. for any $|c|=1$.

  Suppose now that $z\neq0$, and choose
  $\alpha=\overline z/|z|$. Then $|\alpha|=1$ and
  $\alpha z=|z|$. Complex linearity and the pointwise inequality
  $\operatorname{Re}(\alpha f)\leq|\alpha f|=|f|$ give
  \[
    |z|=\operatorname{Re}\left(\alpha\int_X f\,d\mu\right)
       =\int_X\operatorname{Re}(\alpha f)\,d\mu
       \leq\int_X|f|\,d\mu.
  \]
  Equality holds precisely when the nonnegative integrable function
  $|f|-\operatorname{Re}(\alpha f)$ has integral zero, hence when
  $\operatorname{Re}(\alpha f)=|f|$ a.e. A complex number $w$ satisfies
  $\operatorname{Re}w=|w|$ exactly when it is a nonnegative real
  number. Thus equality forces $\alpha f=|f|$ a.e., or
  $f=\overline\alpha|f|$ a.e.

  Conversely, if $f=c|f|$ a.e. for some $|c|=1$, then
  $\int f=c\int|f|$, so taking absolute values gives equality.
  Here invariance of complex integrals under a.e. equality follows
  by applying Corollary~\ref{cor:ae-integrals-sep-22} to the positive
  and negative parts of their real and imaginary components.
\end{proof}

For real-valued functions, equality means that $f\geq0$ a.e. or
$f\leq0$ a.e. For complex-valued functions, all nonzero values must
have the same phase, outside a null set.

\begin{example}
  Take $X=\{0,1\}$ with counting measure, and set $f(0)=a$, $f(1)=b$.
  The integral triangle inequality becomes
  \[
    |a+b|\leq|a|+|b|.
  \]
  Its equality case says that the two nonzero vectors point in the
  same direction.
\end{example}

\begin{figure}[htbp]
  \centering
  \begin{tikzpicture}[>=Latex,font=\small,scale=0.95]
    \begin{scope}
      \draw[->,gray] (-0.2,0)--(4.1,0);
      \draw[->,gray] (0,-0.2)--(0,2.5);
      \draw[->,very thick,blue!75!black] (0,0)--(2.5,0.5)
        node[midway,below] {$a$};
      \draw[->,very thick,orange!80!black] (2.5,0.5)--(3,2)
        node[midway,right] {$b$};
      \draw[->,thick] (0,0)--(3,2) node[midway,above left] {$a+b$};
      \node at (2,-0.65) {$|a+b|<|a|+|b|$};
    \end{scope}
    \begin{scope}[shift={(6,0)}]
      \draw[->,gray] (-0.2,0)--(4.1,0);
      \draw[->,gray] (0,-0.2)--(0,2.5);
      \draw[->,very thick,blue!75!black] (0,0)--(2,1)
        node[midway,above left] {$a$};
      \draw[->,very thick,orange!80!black] (2,1)--(3.5,1.75)
        node[midway,above left] {$b$};
      \node at (2,-0.65) {$|a+b|=|a|+|b|$};
    \end{scope}
  \end{tikzpicture}
  \caption{Different complex phases cause cancellation. Equality in
    the integral triangle inequality requires a common phase almost
    everywhere, as in the aligned vectors on the right.}
  \label{fig:complex-triangle-sep-24}
\end{figure}

\subsection{What integrals determine about a function}

A single integral cannot determine a function, since positive and
negative contributions can cancel. Knowing the integral over every
measurable set determines an integrable function up to a.e. equality.

\begin{proposition}\label{prop:ae-characterization-sep-24}
  For $f,g\in\mathcal L^1(X,\mu)$, the following are equivalent:
  \begin{enumerate}[label=(\roman*)]
    \item $f=g$ a.e.;
    \item $\displaystyle\int_E f\,d\mu=\int_E g\,d\mu$ for every
      $E\in\mathcal M$;
    \item $\displaystyle\int_X|f-g|\,d\mu=0$.
  \end{enumerate}
\end{proposition}

\begin{proof}
  The equivalence of (i) and (iii) is
  Lemma~\ref{lem:zero-integral-sep-22} applied to $|f-g|$. If (iii)
  holds, then for every $E\in\mathcal M$,
  \[
    \left|\int_E f\,d\mu-\int_E g\,d\mu\right|
      \leq\int_E|f-g|\,d\mu
      \leq\int_X|f-g|\,d\mu=0,
  \]
  proving (ii).

  Finally, suppose (ii) holds, and write $f-g=u+iv$. Taking real and
  imaginary parts shows that $\int_E u=\int_E v=0$ for every $E$.
  For $A_n=\{u>1/n\}$,
  \[
    0=\int_{A_n}u\,d\mu\geq\frac1n\mu(A_n),
  \]
  so $\mu(A_n)=0$. Taking the union over $n$ proves
  $\mu(\{u>0\})=0$. Apply the same argument to $-u$, $v$, and $-v$.
  Thus $u=v=0$ a.e., giving (i).
\end{proof}

\begin{remark}[Measurability after a null modification]
  A measurable modification on a null set leaves every integral
  unchanged. If the measure space is complete, arbitrary changes on
  a measurable null set preserve measurability. Without completeness,
  an arbitrary change need not be measurable; one can instead use a
  measurable extension, such as assigning zero on the exceptional
  set, or work with the completion of the measure.
\end{remark}

\subsection{The space \texorpdfstring{$L^1$}{L1}}

If $E\in\mathcal M$, $\mu(X\setminus E)=0$, and $f\colon E\to\bb C$
is measurable for the restricted $\sigma$-algebra, define its extension
\[
  \widetilde f(x)=
  \begin{cases}
    f(x),&x\in E,\\
    0,&x\notin E.
  \end{cases}
\]
This extension is measurable. Define integrability and the integral of
$f$ using $\widetilde f$. Any other measurable extension gives the
same integral, by Proposition~\ref{prop:ae-characterization-sep-24}.
Thus a function need only be defined on a measurable set of full
measure. Similarly, an extended-real-valued measurable function that
is finite a.e. can be changed to zero where it is infinite. It is
integrable precisely when the integral of its absolute value is finite.

\begin{definition}
  For $f,g\in\mathcal L^1(X,\mu)$, write $f\sim g$ if $f=g$ a.e.
  Define
  \[
    L^1(X,\mu)=\mathcal L^1(X,\mu)/{\sim}.
  \]
  An element $[f]\in L^1(X,\mu)$ is an equivalence class of integrable
  functions. Equivalently, one can allow representatives defined only
  on measurable sets of full measure, comparing them on their common
  domain. Set
  \[
    \int_X[f]\,d\mu:=\int_X f\,d\mu.
  \]
  This is well defined by Proposition~\ref{prop:ae-characterization-sep-24}.
\end{definition}

The relation $\sim$ is transitive because a finite union of null sets
is null; reflexivity and symmetry are immediate. Vector operations
$[f]+[g]=[f+g]$ and $c[f]=[cf]$ are independent of representatives for
the same reason. We usually write $f\in L^1$ instead of $[f]\in L^1$,
with the equivalence class understood.

\begin{example}[Point evaluation is not generally well defined]
  For Lebesgue measure,
  \[
    [\mathbf 1_{\bb Q}]=[0]\quad\text{in }L^1(\bb R,m),
  \]
  although $\mathbf 1_{\bb Q}(0)=1$ and the zero function takes value
  zero there. Thus an $L^1(\bb R,m)$ element has no intrinsic value
  at an individual point. This depends on the measure: with counting
  measure there are no nonempty null sets, so a.e. equality is
  pointwise equality.
\end{example}

\begin{proposition}[Completion does not change the quotient space]
  If $(X,\overline{\mathcal M},\overline\mu)$ is the completion of
  $(X,\mathcal M,\mu)$, then the natural map
  \[
    L^1(X,\mu)\longrightarrow L^1(X,\overline\mu),
    \qquad [f]\longmapsto[f],
  \]
  is a bijection preserving integrals and integrals of absolute values.
\end{proposition}

\begin{proof}
  An $\mathcal M$-measurable nonnegative function has the same integral
  for $\mu$ and $\overline\mu$: approximate it increasingly by
  $\mathcal M$-simple functions and apply monotone convergence to
  both measures. The corresponding assertion for integrable
  functions follows componentwise. Also, an $\mathcal M$-measurable
  set is $\overline\mu$-null exactly when it is $\mu$-null, proving
  injectivity of the displayed map.

  For surjectivity, let $f$ be a complex-valued
  $\overline{\mathcal M}$-measurable integrable function. Choose
  $\overline{\mathcal M}$-simple functions $\phi_n\to f$ pointwise.
  By the description of the completion, each measurable set appearing
  in each $\phi_n$ differs from an $\mathcal M$-measurable set by a
  subset of an $\mathcal M$-measurable null set. Replace those sets to
  obtain $\mathcal M$-simple functions $\psi_n$ agreeing with $\phi_n$
  outside such null sets. There is one $N\in\mathcal M$, $\mu(N)=0$,
  containing all these exceptions, since only countably many sets
  were used. The functions $\psi_n\mathbf 1_{X\setminus N}$ converge
  everywhere: to $f$ outside $N$ and to zero on $N$. Their limit $g$
  is $\mathcal M$-measurable and equals $f$ $\overline\mu$-a.e.
  Thus $g$ is integrable and represents the same class as $f$.
\end{proof}

\begin{remark}
  The spaces of actual measurable functions can differ before taking
  the quotient. For example, a non-Borel subset $A$ of the Cantor set
  is Lebesgue measurable and null, so $\mathbf 1_A$ is Lebesgue
  measurable but not Borel measurable. Such subsets exist because
  the Cantor set has continuum cardinality, whereas there are only
  continuum many Borel sets. Nevertheless $\mathbf 1_A$ represents
  the zero class in $L^1$.
\end{remark}

\begin{proposition}[The metric on \texorpdfstring{$L^1$}{L1}]
  \label{prop:l1-metric-sep-24}
  The formula
  \[
    d([f],[g])=\int_X|f-g|\,d\mu
  \]
  defines a metric on $L^1(X,\mu)$. It is induced by the norm
  \[
    \|[f]\|_1=\int_X|f|\,d\mu.
  \]
\end{proposition}

\begin{proof}
  Replacing either representative changes $|f-g|$ only on a null set,
  so $d$ is well defined. It is finite because
  $|f-g|\leq|f|+|g|$, and it is nonnegative and symmetric. Moreover,
  \[
    d([f],[g])=0\quad\Longleftrightarrow\quad f=g\text{ a.e.}
      \quad\Longleftrightarrow\quad[f]=[g].
  \]
  The pointwise inequality $|f-h|\leq|f-g|+|g-h|$ gives
  \[
    d([f],[h])\leq d([f],[g])+d([g],[h]).
  \]
  Finally, $\|c[f]\|_1=|c|\|[f]\|_1$, and the remaining norm axioms
  follow from the same arguments.
\end{proof}

On actual functions, the same formula is only a pseudometric: distinct
functions can have distance zero. Passing to equivalence classes is
exactly what makes the separation axiom valid.

\FloatBarrier

\subsection{The dominated convergence theorem}

Monotone convergence requires an increasing nonnegative sequence.
Fatou's lemma removes monotonicity but gives only one inequality.
An integrable bound permits passage to the limit even for complex
functions and gives convergence in the $L^1$ metric.

\begin{theorem}[Dominated convergence theorem]
  \label{thm:dct-sep-24}
  Let $f_n\colon X\to\bb C$ be measurable, and suppose $f_n\to f$
  a.e., where $f$ is measurable. Suppose there is a nonnegative
  $g\in\mathcal L^1(X,\mu)$ such that $|f_n|\leq g$ a.e. for every
  $n$. Then $f_n,f\in\mathcal L^1(X,\mu)$ and
  \[
    \lim_{n\to\infty}\int_X|f_n-f|\,d\mu=0.
  \]
  Consequently,
  \[
    \lim_{n\to\infty}\int_X f_n\,d\mu=\int_X f\,d\mu.
  \]
\end{theorem}

\begin{proof}
  Take a common measurable null set outside which convergence and
  all the bounds hold, and redefine the functions to be zero there.
  This does not change their integrals. Thus we may assume
  $f_n\to f$ and $|f_n|\leq g$ everywhere. Taking limits gives
  $|f|\leq g$, so $f$ and all $f_n$ are integrable.

  Put $h_n=|f_n-f|$. Then $h_n\to0$ and $0\leq h_n\leq2g$.
  Apply Fatou's lemma to the nonnegative functions $2g-h_n$:
  \begin{align*}
    2\int_X g\,d\mu
      &=\int_X\liminf_n(2g-h_n)\,d\mu\\
      &\leq\liminf_n\int_X(2g-h_n)\,d\mu\\
      &=2\int_X g\,d\mu-\limsup_n\int_X h_n\,d\mu.
  \end{align*}
  All integrals here are finite, and
  $0\leq\int h_n\leq2\int g$. Subtracting the finite quantity
  $2\int g$ yields $\limsup_n\int h_n\leq0$. Nonnegativity now
  implies $\int h_n\to0$. Finally,
  \[
    \left|\int_X f_n\,d\mu-\int_X f\,d\mu\right|
      \leq\int_X|f_n-f|\,d\mu\longrightarrow0.
  \]
\end{proof}

\begin{remark}
  If the pointwise limit is initially defined only on a measurable
  set of full measure, extend it by zero on the complement. It is
  then measurable as a pointwise limit of the similarly modified
  $f_n$. Thus the theorem also applies naturally to a.e. defined
  representatives of $L^1$ elements.
\end{remark}

\begin{example}[Why the dominating function must be integrable]
  \label{ex:dominating-envelope-sep-24}
  The smallest pointwise nonnegative bound for a sequence is
  $g_0=\sup_n|f_n|$. Consider the following sequences on $\bb R$.
  \begin{enumerate}[label=(\roman*)]
    \item For $f_n=\mathbf 1_{(n,n+1]}$, $n\geq1$, we have
      $f_n\to0$ pointwise and $\int f_n\,dm=1$. But
      \[
        g_0=\mathbf 1_{(1,+\infty)},\qquad\int g_0\,dm=+\infty.
      \]
      Thus there is no integrable dominating function.
    \item Let $f_n=n^{-1}\mathbf 1_{(0,n]}$. Then
      $\sup_x|f_n(x)|=1/n\to0$, so the convergence to zero is even
      uniform, while $\int f_n\,dm=1$ for every $n$. Here
      \[
        g_0(x)=
        \begin{cases}
          0,&x\leq0,\\
          1/k,&k-1<x\leq k,\quad k\in\bb N.
        \end{cases}
      \]
      Indeed, for $k-1<x\leq k$, the nonzero terms are those with
      $n\geq k$, and their largest value is $1/k$. Therefore
      \[
        \int_{\bb R}g_0\,dm=\sum_{k=1}^{\infty}\frac1k=+\infty.
      \]
      Uniform convergence on a space of infinite measure does not
      by itself imply convergence of integrals.
  \end{enumerate}
  The same obstruction holds for a.e. domination: after taking the
  union of countably many exceptional null sets, any common bound
  dominates $g_0$ a.e., and hence also has infinite integral.
\end{example}

\begin{figure}[!htbp]
  \centering
  \begin{tikzpicture}[x=1.3cm,y=2.6cm,>=Latex,font=\small]
    \foreach \k in {1,...,6} {
      \fill[blue!12] (\k-1,0) rectangle (\k,1/\k);
      \draw[blue!75!black,very thick] (\k-1,1/\k)--(\k,1/\k);
      \draw[blue!40,dotted] (\k,0)--(\k,1/\k);
    }
    \draw[->] (-0.1,0)--(6.65,0) node[right] {$x$};
    \draw[->] (0,-0.04)--(0,1.2) node[above] {value};
    \draw[orange!85!black,thick,dashed] (0,0.333333)--(3,0.333333)
      --(3,0) node[below=15pt] {$f_3=\frac13\mathbf 1_{(0,3]}$};
    \foreach \k in {1,...,6} {
      \draw (\k,0.015)--(\k,-0.015) node[below] {$\k$};
    }
    \node[left] at (0,1) {$1$};
    \node[blue!75!black,right] at (1.1,0.92) {$g_0=\sup_n f_n$};
    \node[blue!75!black] at (4.5,0.75)
      {$\displaystyle\int g_0\,dm=\sum_{k=1}^{\infty}\frac1k=+\infty$};
  \end{tikzpicture}
  \caption{The envelope of $f_n=n^{-1}\mathbf 1_{(0,n]}$ has a
    divergent harmonic tail. Every $f_n$ is integrable, but no single
    integrable function dominates the sequence.}
  \label{fig:dominating-envelope-sep-24}
\end{figure}

\FloatBarrier
